\section{Phase 6 Derivation Draft: Qualified Contract-Manufacturing Market and Labor-Market Clearing}

This phase opens the market-clearing block of the benchmark MAH model. The task here is deliberately narrower than a proof of stationary competitive equilibrium. We take as given the optimized static controls derived in Phase 2, in particular $x_B^*(z;p_m,M)$ and $m^*(z;p_m,w)$; the fixed-price value-function block completed in Phase 3; the entrant masses $n_E^A(M),n_E^B(M),n_E^C(M)$ from Phase 5; the compact state space $Z$; and the earlier regularity conditions already imposed on those benchmark objects. We also take as given candidate incumbent measures $\mu_A,\mu_B,\mu_C$. At this stage these measures are treated only as candidate finite Borel measures, or equivalently candidate incumbent measures. They are not yet treated as stationary objects, because invariant-measure arguments belong to the next phase rather than to the opening setup of Phase 6.

\subsection{Benchmark Fidelity Check}

\paragraph{Step 1: exact benchmark objects relevant for Phase 6.}
The objects that matter in the opening setup of Phase 6 are all inherited from the benchmark manuscript and from the logic already completed in Phases 1--5. From Phase 2 we import the optimized operating controls $x_B^*(z;p_m,M)$ and $m^*(z;p_m,w)$, together with the optimized payoff objects $\Pi_A(z;w,M)$, $\Pi_B(z;p_m,M)$, and $\Pi_C(z;p_m,w,M)$ as background objects already derived and not to be rewritten into a generic $\Pi_j$ notation. From Phase 3 we inherit the fixed-price value-function block, but only as an already-completed dynamic input and not as a new object to be reproved here. From Phase 5 we inherit the entrant masses $n_E^A(M),n_E^B(M),n_E^C(M)$. In addition, this phase takes as given the compact state space $Z$, the candidate incumbent measures $\mu_A,\mu_B,\mu_C$ understood as candidate finite Borel measures, the reduced-form commercialization-success object $\lambda_B(z,p_m,M)$, the type-specific labor objects $\ell_A(z;w)$, $\ell_B(z;p_m,M)$, and $\ell_C(z;p_m,w)$, the entrant setup labor coefficients $\ell_E^j$, and the exogenous labor supply $\bar L$.

\paragraph{Step 2: exact benchmark market objects and normalization.}
The contract-manufacturing block in the benchmark uses a specific normalization that must be stated explicitly here: one successfully commercialized type-$B$ original-drug idea requires one unit of qualified external manufacturing service. This is a benchmark normalization embedded in the market-clearing block. It is not a new implementation problem, not a separate contracting layer, and not a new primitive optimization block. Under that normalization, qualified contract-manufacturing demand and supply are
\begin{align}
D_m(p_m,M)
&=\int_Z x_B^*(z;p_m,M)\,\lambda_B(z,p_m,M)\,\mu_B(dz),
\label{eq:p6_Dm}\\
S_m(p_m,w)
&=\int_Z m^*(z;p_m,w)\,\mu_C(dz),
\label{eq:p6_Sm}
\end{align}
and market clearing requires
\begin{equation}
D_m(p_m,M)=S_m(p_m,w).
\label{eq:p6_pm_clear}
\end{equation}

The benchmark labor-demand object is written directly as
\begin{equation}
\begin{aligned}
L_D(M)={}&\int_Z \ell_A(z;w)\,\mu_A(dz)
+\int_Z \ell_B(z;p_m,M)\,\mu_B(dz)\\
&+\int_Z \ell_C(z;p_m,w)\,\mu_C(dz)
+\sum_{j\in\{A,B,C\}}\ell_E^j n_E^j(M),
\end{aligned}
\end{equation}
and labor-market clearing in the benchmark is
\begin{equation}
L_D(M)=\bar L.
\end{equation}

\paragraph{Step 3: what this phase can complete faithfully, and what remains pending.}
At the opening stage, Phase 6 can faithfully restate the benchmark market-clearing objects conditional on candidate prices, candidate incumbent measures, and entrant masses. It can also prepare the later derivations by making explicit which already-derived policy objects enter those market expressions. What remains pending is equally important: this phase-opening setup does not establish invariant measures, does not establish stationary competitive-equilibrium existence, and does not solve any joint price-measure fixed-point problem. Those claims belong to later phases and should not be silently imported here.

\paragraph{Step 4: why Phase 6 must be written before Phases 7 and 8.}
Phase 6 is needed because the later stationary-distribution and equilibrium blocks require a precise market-clearing interface. Phase 7 cannot formulate stationary consistency in an economically meaningful way unless the contract-manufacturing and labor-market objects have already been written down in their benchmark form. Phase 8 cannot close the full stationary competitive equilibrium until the value-function block, the entry block, the market-clearing block, and the stationary-distribution block have all been stated separately and then combined. For that reason, Phase 6 isolates the benchmark market objects first and records them in a form that can be passed directly to the next two phases.

The next parts of Phase 6 therefore derive the contract-manufacturing market object and the labor-market object conditional on candidate prices, candidate incumbent measures, and entrant masses.

\subsection{Part 1: Qualified Contract-Manufacturing Market}

\paragraph{Benchmark object restatement.}
In the benchmark manuscript, the qualified contract-manufacturing market is summarized by the aggregate demand object
\[
D_m(p_m,M)=\int_Z x_B^*(z;p_m,M)\,\lambda_B(z,p_m,M)\,\mu_B(dz),
\]
the aggregate supply object
\[
S_m(p_m,w)=\int_Z m^*(z;p_m,w)\,\mu_C(dz),
\]
and the market-clearing restriction
\[
D_m(p_m,M)=S_m(p_m,w).
\]
These are the exact benchmark objects and they are kept here without redesign. The measures $\mu_B$ and $\mu_C$ are treated in this phase as candidate finite Borel measures. They are not yet proved to be invariant stationary measures. The controls $x_B^*(z;p_m,M)$ and $m^*(z;p_m,w)$ are already-solved policy objects imported from Phase 2 and are not re-optimized in this section.

\paragraph{Economic role.}
The benchmark normalization is that one successfully commercialized type-$B$ original-drug idea requires one unit of qualified external manufacturing service. This normalization explains why the product $x_B^*(z;p_m,M)\lambda_B(z,p_m,M)$ enters the demand object directly. For a type-$B$ firm in state $z$, the control $x_B^*(z;p_m,M)$ is the optimized original-drug idea-generation intensity, and $\lambda_B(z,p_m,M)$ is the reduced-form success rate with which those ideas are commercialized through qualified external manufacturing. Their product therefore measures the expected amount of qualified external manufacturing services demanded by that firm-state. Integrating this object against the candidate type-$B$ measure $\mu_B$ produces the expected industry demand for qualified external manufacturing generated by type-$B$ firms.

The supply object is interpreted in the parallel way on the type-$C$ side. For a type-$C$ firm in state $z$, the policy $m^*(z;p_m,w)$ is the optimized quantity of qualified contract-manufacturing services supplied at candidate prices $(p_m,w)$. Integrating this supply policy against the candidate type-$C$ measure $\mu_C$ produces the industry supply of qualified external manufacturing generated by type-$C$ firms. Thus the contract-manufacturing market compares, on the demand side, the expected outsourcing needs created by type-$B$ commercialization and, on the supply side, the production capacity offered by type-$C$ incumbents.

\paragraph{Mathematical problem.}
Fix a regime $M\in\{0,1\}$, a candidate qualified-manufacturing price $p_m$, a candidate wage $w$, and candidate finite Borel measures $\mu_B$ and $\mu_C$. The mathematical task in this part is limited to three questions. First, are the two integrals defining $D_m(p_m,M)$ and $S_m(p_m,w)$ well defined and finite? Second, under local compactness and continuity conditions in prices, do these aggregate objects vary continuously with prices? Third, once the two aggregate objects are constructed, how should the benchmark market-clearing restriction be recorded as an excess-demand object for later equilibrium analysis? This part does not solve the restriction, does not claim that a clearing price exists, and does not claim that a clearing price is unique.

\paragraph{Supplemental assumptions.}
\begin{assumption}[Measurability and finite candidate measures]
\label{ass:p6_cm_meas}
For each fixed triple $(p_m,w,M)$, the map
\[
z\mapsto x_B^*(z;p_m,M)\lambda_B(z,p_m,M)
\]
is Borel measurable on $Z$, and the map
\[
z\mapsto m^*(z;p_m,w)
\]
is Borel measurable on $Z$. The candidate measures $\mu_B$ and $\mu_C$ are finite Borel measures on $Z$.
\end{assumption}

\begin{assumption}[Boundedness at fixed candidate prices]
\label{ass:p6_cm_bound}
For each fixed triple $(p_m,w,M)$, there exist finite constants $B_D(p_m,M)$ and $B_S(p_m,w)$ such that
\[
\left|x_B^*(z;p_m,M)\lambda_B(z,p_m,M)\right|\le B_D(p_m,M)
\quad\text{for all } z\in Z,
\]
and
\[
\left|m^*(z;p_m,w)\right|\le B_S(p_m,w)
\quad\text{for all } z\in Z.
\]
\end{assumption}

\begin{assumption}[Joint continuity on compact local price domains]
\label{ass:p6_cm_cont}
Fix a regime $M$. Let $\mathcal P_m^{\mathrm{loc}}\subset\mathbb R_+$ be a compact local domain for the qualified-manufacturing price and let $\mathcal W^{\mathrm{loc}}\subset\mathbb R_+$ be a compact local domain for the wage. The map
\[
(z,p_m)\mapsto x_B^*(z;p_m,M)\lambda_B(z,p_m,M)
\]
is jointly continuous on $Z\times \mathcal P_m^{\mathrm{loc}}$, and the map
\[
(z,p_m,w)\mapsto m^*(z;p_m,w)
\]
is jointly continuous on $Z\times \mathcal P_m^{\mathrm{loc}}\times \mathcal W^{\mathrm{loc}}$.
\end{assumption}

\begin{proposition}[Well-definedness and finiteness]
Fix a regime $M$, candidate prices $(p_m,w)$, and candidate finite Borel measures $(\mu_B,\mu_C)$. Under Assumptions~\ref{ass:p6_cm_meas} and~\ref{ass:p6_cm_bound}, the objects $D_m(p_m,M)$ and $S_m(p_m,w)$ are well defined as finite real numbers.
\end{proposition}

\begin{proof}
Fix $(p_m,w,M,\mu_B,\mu_C)$. Define the two integrands
\[
f_D(z;p_m,M):=x_B^*(z;p_m,M)\lambda_B(z,p_m,M)
\]
and
\[
f_S(z;p_m,w):=m^*(z;p_m,w).
\]
By Assumption~\ref{ass:p6_cm_meas}, both $f_D(\cdot;p_m,M)$ and $f_S(\cdot;p_m,w)$ are Borel measurable on $Z$. Therefore each is an admissible integrand for integration with respect to the corresponding candidate Borel measure.

Next, use Assumption~\ref{ass:p6_cm_bound}. For every $z\in Z$,
\[
|f_D(z;p_m,M)|=\left|x_B^*(z;p_m,M)\lambda_B(z,p_m,M)\right|\le B_D(p_m,M).
\]
Integrating this inequality with respect to $\mu_B$ gives
\[
\int_Z |f_D(z;p_m,M)|\,\mu_B(dz)
\le
\int_Z B_D(p_m,M)\,\mu_B(dz).
\]
Because $B_D(p_m,M)$ is a constant in $z$, the right-hand side simplifies to
\[
\int_Z B_D(p_m,M)\,\mu_B(dz)
=
B_D(p_m,M)\mu_B(Z).
\]
Assumption~\ref{ass:p6_cm_meas} states that $\mu_B$ is a finite Borel measure, so $\mu_B(Z)<\infty$. Assumption~\ref{ass:p6_cm_bound} states that $B_D(p_m,M)<\infty$. Hence
\[
\int_Z |f_D(z;p_m,M)|\,\mu_B(dz)<\infty.
\]
Therefore $f_D(\cdot;p_m,M)$ is absolutely integrable under $\mu_B$, and the integral defining $D_m(p_m,M)$ exists and is finite.

Apply the same argument to supply. For every $z\in Z$,
\[
|f_S(z;p_m,w)|=|m^*(z;p_m,w)|\le B_S(p_m,w).
\]
Integrating with respect to $\mu_C$ yields
\[
\int_Z |f_S(z;p_m,w)|\,\mu_C(dz)
\le
\int_Z B_S(p_m,w)\,\mu_C(dz)
=
B_S(p_m,w)\mu_C(Z).
\]
Again, $\mu_C(Z)<\infty$ by Assumption~\ref{ass:p6_cm_meas} and $B_S(p_m,w)<\infty$ by Assumption~\ref{ass:p6_cm_bound}. Hence
\[
\int_Z |f_S(z;p_m,w)|\,\mu_C(dz)<\infty.
\]
Therefore the integral defining $S_m(p_m,w)$ also exists and is finite.

We have now shown, step by step, that both benchmark market objects are well defined as finite real numbers for fixed candidate prices and fixed candidate measures.
\end{proof}

\begin{proposition}[Continuity on compact local price domains]
Fix a regime $M$ and candidate finite Borel measures $(\mu_B,\mu_C)$. Under Assumptions~\ref{ass:p6_cm_meas}--\ref{ass:p6_cm_cont}, the demand object $D_m(\cdot,M)$ is continuous on $\mathcal P_m^{\mathrm{loc}}$, and the supply object $S_m(\cdot,\cdot)$ is jointly continuous on $\mathcal P_m^{\mathrm{loc}}\times \mathcal W^{\mathrm{loc}}$. In particular, for each fixed wage $w\in\mathcal W^{\mathrm{loc}}$, the map $p_m\mapsto S_m(p_m,w)$ is continuous on $\mathcal P_m^{\mathrm{loc}}$.
\end{proposition}

\begin{proof}
We begin with demand. Define
\[
g_D(z,p_m):=x_B^*(z;p_m,M)\lambda_B(z,p_m,M).
\]
By Assumption~\ref{ass:p6_cm_cont}, the map $g_D$ is jointly continuous on the compact set $Z\times\mathcal P_m^{\mathrm{loc}}$. A continuous function on a compact set is bounded by the Extreme Value Theorem. Therefore there exists a finite constant $K_D$ such that
\[
|g_D(z,p_m)|\le K_D
\quad\text{for every } (z,p_m)\in Z\times\mathcal P_m^{\mathrm{loc}}.
\]

Now let $\{p_m^n\}_{n\ge 1}$ be any sequence in $\mathcal P_m^{\mathrm{loc}}$ such that $p_m^n\to p_m$ for some $p_m\in\mathcal P_m^{\mathrm{loc}}$. Fix any state $z\in Z$. Joint continuity of $g_D$ implies
\[
g_D(z,p_m^n)\to g_D(z,p_m)
\quad\text{as } n\to\infty.
\]
Moreover, for every $n$ and every $z\in Z$,
\[
|g_D(z,p_m^n)|\le K_D.
\]
Because $\mu_B$ is a finite measure, the constant function $z\mapsto K_D$ is $\mu_B$-integrable:
\[
\int_Z K_D\,\mu_B(dz)=K_D\mu_B(Z)<\infty.
\]
Therefore the Dominated Convergence Theorem applies. It follows that
\[
\lim_{n\to\infty} D_m(p_m^n,M)
=
\lim_{n\to\infty}\int_Z g_D(z,p_m^n)\,\mu_B(dz)
=
\int_Z g_D(z,p_m)\,\mu_B(dz)
=
D_m(p_m,M).
\]
This proves continuity of $D_m(\cdot,M)$ on $\mathcal P_m^{\mathrm{loc}}$.

We now turn to supply. Define
\[
g_S(z,p_m,w):=m^*(z;p_m,w).
\]
By Assumption~\ref{ass:p6_cm_cont}, the map $g_S$ is jointly continuous on the compact set $Z\times\mathcal P_m^{\mathrm{loc}}\times\mathcal W^{\mathrm{loc}}$. Again by the Extreme Value Theorem, there exists a finite constant $K_S$ such that
\[
|g_S(z,p_m,w)|\le K_S
\quad\text{for every } (z,p_m,w)\in Z\times\mathcal P_m^{\mathrm{loc}}\times\mathcal W^{\mathrm{loc}}.
\]

Let $\{(p_m^n,w^n)\}_{n\ge 1}$ be any sequence in $\mathcal P_m^{\mathrm{loc}}\times\mathcal W^{\mathrm{loc}}$ such that
\[
(p_m^n,w^n)\to(p_m,w)
\]
for some $(p_m,w)\in\mathcal P_m^{\mathrm{loc}}\times\mathcal W^{\mathrm{loc}}$. Fix any state $z\in Z$. Joint continuity of $g_S$ implies
\[
g_S(z,p_m^n,w^n)\to g_S(z,p_m,w)
\quad\text{as } n\to\infty.
\]
At the same time, for every $n$ and every $z\in Z$,
\[
|g_S(z,p_m^n,w^n)|\le K_S.
\]
Since $\mu_C$ is finite,
\[
\int_Z K_S\,\mu_C(dz)=K_S\mu_C(Z)<\infty.
\]
The Dominated Convergence Theorem therefore yields
\[
\lim_{n\to\infty} S_m(p_m^n,w^n)
=
\lim_{n\to\infty}\int_Z g_S(z,p_m^n,w^n)\,\mu_C(dz)
=
\int_Z g_S(z,p_m,w)\,\mu_C(dz)
=
S_m(p_m,w).
\]
Thus $S_m(\cdot,\cdot)$ is jointly continuous on $\mathcal P_m^{\mathrm{loc}}\times\mathcal W^{\mathrm{loc}}$. The continuity of $p_m\mapsto S_m(p_m,w)$ for each fixed $w$ is the special case in which only the first price argument varies.
\end{proof}

\paragraph{Definition of market-clearing restriction and excess demand.}
Let $\mu:=(\mu_B,\mu_C)$ denote the candidate measure pair entering the contract-manufacturing market. The benchmark market-clearing restriction can be written as
\[
D_m(p_m,M)-S_m(p_m,w)=0.
\]
Accordingly, define the contract-manufacturing excess-demand object by
\[
Z_m(p_m,w,M;\mu):=D_m(p_m,M)-S_m(p_m,w).
\]
Under Assumptions~\ref{ass:p6_cm_meas} and~\ref{ass:p6_cm_bound}, this object is well defined for fixed candidate prices and fixed candidate measures. Under Assumption~\ref{ass:p6_cm_cont}, it is continuous on compact local price domains. At the present stage of Phase 6, it serves only to record the benchmark market-clearing restriction in a form suitable for later equilibrium analysis.

The next part derives the labor-market object under the same conditional logic.

\subsection{Part 2: Labor Market}

\paragraph{Benchmark object restatement.}
The benchmark labor-market block is summarized by the aggregate labor-demand object
\[
\begin{aligned}
L_D(M)={}&\int_Z \ell_A(z;w)\,\mu_A(dz)
+\int_Z \ell_B(z;p_m,M)\,\mu_B(dz)\\
&+\int_Z \ell_C(z;p_m,w)\,\mu_C(dz)
+\sum_{j\in\{A,B,C\}}\ell_E^j n_E^j(M),
\end{aligned}
\]
and the labor-market-clearing restriction
\[
L_D(M)=\bar L.
\]
We keep this benchmark object exactly as written. In the present phase it is evaluated at candidate prices $(p_m,w)$, candidate finite Borel measures $(\mu_A,\mu_B,\mu_C)$, and entrant masses inherited from Phase 5, but we do not replace the benchmark notation with a new public object.

\paragraph{Economic role.}
Each term in $L_D(M)$ has a direct economic interpretation. The first term is incumbent type-$A$ labor demand integrated across the candidate measure of integrated firms. The second term is incumbent type-$B$ labor demand integrated across the candidate measure of research-specialized firms. The third term is incumbent type-$C$ labor demand integrated across the candidate measure of manufacturing-specialized firms. The last term is entrant setup labor: each entrant of type $j\in\{A,B,C\}$ uses $\ell_E^j$ units of setup labor, and the total entrant labor contribution is therefore $\ell_E^j n_E^j(M)$ aggregated across entrant types.

The key modeling point is that $\ell_A(z;w)$, $\ell_B(z;p_m,M)$, and $\ell_C(z;p_m,w)$ are not derived here from a new capital-labor choice problem. They are labor-demand objects recovered from the already optimized flow-payoff blocks. In other words, the labor market takes the operating problems from earlier phases as given and asks how much labor those optimized firm decisions imply in the aggregate once they are weighted by candidate incumbent measures and augmented by entrant setup labor from the Phase 5 entry block.

\paragraph{Mathematical problem.}
Fix a regime $M\in\{0,1\}$, candidate prices $(p_m,w)$, candidate finite Borel measures $(\mu_A,\mu_B,\mu_C)$, and entrant masses $\bigl(n_E^A(M),n_E^B(M),n_E^C(M)\bigr)$. The purpose of this part is modest. We want to show that the benchmark object $L_D(M)$ is well defined for such candidate inputs, and that under mild local regularity it varies continuously with prices on compact local price domains. Once that has been established, we can record the labor-market-clearing restriction in excess-demand form for later equilibrium analysis. We do not solve for a clearing wage here, and we do not claim wage uniqueness.

\paragraph{Minimal regularity.}
To make the benchmark labor object usable, we need only a small amount of additional structure. First, for fixed candidate prices, the labor-demand objects $\ell_A(\cdot;w)$, $\ell_B(\cdot;p_m,M)$, and $\ell_C(\cdot;p_m,w)$ should be Borel measurable in $z$, because each must be integrable against the corresponding candidate measure. Second, for fixed candidate prices, each labor-demand object should be bounded on the compact state space $Z$, so that integration against finite measures produces finite values. Third, the entrant coefficients $\ell_E^j$ and the entrant masses $n_E^j(M)$ must be finite, so the entrant labor term is a finite sum. Finally, if we want continuity in candidate prices on compact local domains, we need the three labor-demand objects to vary continuously with their price arguments on those domains.

For reference in later cross-references, we summarize these conditions as follows.
\begin{assumption}[Measurability and finite labor objects]
\label{ass:p6_lab_meas}
For each fixed triple $(p_m,w,M)$, the maps $z\mapsto \ell_A(z;w)$, $z\mapsto \ell_B(z;p_m,M)$, and $z\mapsto \ell_C(z;p_m,w)$ are Borel measurable on $Z$. The candidate measures $\mu_A,\mu_B,\mu_C$ are finite Borel measures on $Z$. The entrant coefficients $\ell_E^A,\ell_E^B,\ell_E^C$ and entrant masses $n_E^A(M),n_E^B(M),n_E^C(M)$ are finite.
\end{assumption}

\begin{assumption}[Bounded labor demand at fixed candidate prices]
\label{ass:p6_lab_bound}
For each fixed triple $(p_m,w,M)$, there exist finite constants $B_A(w)$, $B_B(p_m,M)$, and $B_C(p_m,w)$ such that
\[
|\ell_A(z;w)|\le B_A(w),\qquad
|\ell_B(z;p_m,M)|\le B_B(p_m,M),\qquad
|\ell_C(z;p_m,w)|\le B_C(p_m,w)
\]
for all $z\in Z$.
\end{assumption}

\begin{assumption}[Local continuity in candidate prices]
\label{ass:p6_lab_cont}
Fix a regime $M$. Let $\mathcal P_m^{\mathrm{loc}}\subset\mathbb R_+$ and $\mathcal W^{\mathrm{loc}}\subset\mathbb R_+$ be compact local price domains. The map $(z,w)\mapsto \ell_A(z;w)$ is jointly continuous on $Z\times\mathcal W^{\mathrm{loc}}$, the map $(z,p_m)\mapsto \ell_B(z;p_m,M)$ is jointly continuous on $Z\times\mathcal P_m^{\mathrm{loc}}$, and the map $(z,p_m,w)\mapsto \ell_C(z;p_m,w)$ is jointly continuous on $Z\times\mathcal P_m^{\mathrm{loc}}\times\mathcal W^{\mathrm{loc}}$.
\end{assumption}

\begin{proposition}[Aggregate labor demand is well defined]
Fix a regime $M$, candidate prices $(p_m,w)$, candidate finite Borel measures $(\mu_A,\mu_B,\mu_C)$, and entrant masses $\bigl(n_E^A(M),n_E^B(M),n_E^C(M)\bigr)$. Under Assumptions~\ref{ass:p6_lab_meas} and~\ref{ass:p6_lab_bound}, the benchmark aggregate labor-demand object $L_D(M)$ is well defined and finite.
\end{proposition}

\begin{proof}
Write $L_D(M)$ as the sum of four components:
\[
L_D(M)=L_A+L_B+L_C+L_E,
\]
where
\[
L_A:=\int_Z \ell_A(z;w)\,\mu_A(dz),\quad
L_B:=\int_Z \ell_B(z;p_m,M)\,\mu_B(dz),\quad
L_C:=\int_Z \ell_C(z;p_m,w)\,\mu_C(dz),
\]
and
\[
L_E:=\sum_{j\in\{A,B,C\}}\ell_E^j n_E^j(M).
\]
By Assumption~\ref{ass:p6_lab_meas}, each incumbent labor-demand object is Borel measurable in $z$, so the three integrals are well defined as measurable integrals. By Assumption~\ref{ass:p6_lab_bound},
\[
|\ell_A(z;w)|\le B_A(w)\quad\text{for all } z\in Z.
\]
Hence
\[
\int_Z |\ell_A(z;w)|\,\mu_A(dz)\le B_A(w)\mu_A(Z)<\infty,
\]
because $\mu_A$ is finite by Assumption~\ref{ass:p6_lab_meas}. The same argument gives
\[
\int_Z |\ell_B(z;p_m,M)|\,\mu_B(dz)\le B_B(p_m,M)\mu_B(Z)<\infty
\]
and
\[
\int_Z |\ell_C(z;p_m,w)|\,\mu_C(dz)\le B_C(p_m,w)\mu_C(Z)<\infty.
\]
So the three incumbent components are finite.

For the entrant term, Assumption~\ref{ass:p6_lab_meas} gives finiteness of each coefficient $\ell_E^j$ and each entrant mass $n_E^j(M)$. Therefore each product $\ell_E^j n_E^j(M)$ is finite, and because there are only three entrant types, the sum $L_E$ is finite. Combining the four parts shows that $L_D(M)$ is well defined and finite.
\end{proof}

\begin{proposition}[Continuity on compact local price domains]
Fix a regime $M$, candidate finite Borel measures $(\mu_A,\mu_B,\mu_C)$, and entrant masses $\bigl(n_E^A(M),n_E^B(M),n_E^C(M)\bigr)$. Under Assumptions~\ref{ass:p6_lab_meas}--\ref{ass:p6_lab_cont}, the benchmark labor-demand object $L_D(M)$ is jointly continuous in $(p_m,w)$ on $\mathcal P_m^{\mathrm{loc}}\times\mathcal W^{\mathrm{loc}}$.
\end{proposition}

\begin{proof}
The entrant term does not vary with candidate prices in this phase, so continuity reduces to the three incumbent integrals. Under Assumption~\ref{ass:p6_lab_cont}, each labor-demand function is jointly continuous on a compact set. By the Extreme Value Theorem, each is therefore uniformly bounded on the relevant compact domain. Because the candidate measures are finite, these uniform bounds are integrable dominating functions. For any convergent sequence of candidate prices $(p_m^n,w^n)\to(p_m,w)$, pointwise convergence of the three labor-demand functions follows from joint continuity, and the Dominated Convergence Theorem then implies convergence of each corresponding integral. Summing the three incumbent terms and the constant entrant term shows that $L_D(M)$ is jointly continuous in $(p_m,w)$ on compact local price domains.
\end{proof}

\paragraph{Labor-market-clearing restriction and labor excess demand.}
The benchmark labor-market-clearing restriction is
\[
L_D(M)=\bar L.
\]
It is convenient for later equilibrium analysis to record the same condition through the labor excess-demand object
\[
Z_\ell(p_m,w,M;\mu,n_E):=L_D(M)-\bar L,
\]
where $\mu:=(\mu_A,\mu_B,\mu_C)$ and $n_E:=\bigl(n_E^A(M),n_E^B(M),n_E^C(M)\bigr)$. Under Assumptions~\ref{ass:p6_lab_meas} and~\ref{ass:p6_lab_bound}, this object is well defined for candidate prices, candidate measures, and entrant masses; under Assumption~\ref{ass:p6_lab_cont}, it is continuous on compact local price domains. At this stage, however, it serves only as a compact labor-market restriction for later equilibrium analysis.

The remaining part of Phase 6 will summarize the conditional market objects for use in later equilibrium analysis.

\subsection{Part 3: Market-Clearing Restrictions, Excess-Demand Notation, and End-of-Phase Recap}

\paragraph{Gathering the market objects.}
Let
\[
\mu:=(\mu_A,\mu_B,\mu_C)
\qquad\text{and}\qquad
n_E:=\bigl(n_E^A(M),n_E^B(M),n_E^C(M)\bigr).
\]
Conditional on candidate prices $(p_m,w)$, candidate incumbent measures $\mu$, and entrant masses $n_E$, Phase 6 has constructed the contract-manufacturing excess-demand object
\[
Z_m(p_m,w,M;\mu):=D_m(p_m,M)-S_m(p_m,w)
\]
and the labor excess-demand object
\[
Z_\ell(p_m,w,M;\mu,n_E):=L_D(M)-\bar L.
\]
The corresponding market-clearing restrictions are
\[
Z_m(p_m,w,M;\mu)=0
\qquad\text{and}\qquad
Z_\ell(p_m,w,M;\mu,n_E)=0.
\]
These are conditional market restrictions. At this stage, Phase 6 has not shown that the same candidate incumbent measures are invariant under optimal behavior, exit, and entrant inflows.

\paragraph{What Phase 6 has established.}
Phase 6 has established that the benchmark contract-manufacturing market objects are defined, that the benchmark aggregate labor-demand object is defined, and that both market-clearing restrictions can be written in excess-demand form. In particular, the objects $Z_m(p_m,w,M;\mu)$ and $Z_\ell(p_m,w,M;\mu,n_E)$ are now available as compact summaries of the market side of the model for later equilibrium analysis. What Phase 6 has not established is equally important. It has not proved an invariant-measure result, it has not proved the existence of a stationary competitive equilibrium, and it has not solved any joint fixed point over value functions, policy behavior, candidate incumbent measures, entrant masses, and prices.

\paragraph{Handoff to later phases.}
Phase 7 must take selected policy behavior, transition kernels, exit, and entrant inflows and combine them into a law of motion for incumbent measures; only then can invariant measures be defined and studied. Phase 8 must then assemble value functions, operating controls, organizational choices, prices, entrant masses, and invariant measures into a stationary competitive-equilibrium definition and an existence framework.
